
\clearpage
\begin{landscape}
\singlespacing
\begin{center}\small\bfseries
ПРИЛОЖЕНИЕ 2 --- Сравнительная таблица характеристик источников энтропии
\end{center}
\nopagebreak\vspace{0.15\baselineskip}\nopagebreak
\centering
{\scriptsize\setlength{\tabcolsep}{2.6pt}\renewcommand{\arraystretch}{0.9}%
{\def\LTcaptype{none} % do not increment counter
\begin{longtable}[]{@{}|%
  >{\raggedright\arraybackslash}p{(\linewidth - 16\tabcolsep - 1pt - 9\arrayrulewidth - 6pt) * \real{0.1251}}|%
  >{\raggedright\arraybackslash}p{(\linewidth - 16\tabcolsep - 1pt - 9\arrayrulewidth - 6pt) * \real{0.1251}}|%
  >{\raggedright\arraybackslash}p{(\linewidth - 16\tabcolsep - 1pt - 9\arrayrulewidth - 6pt) * \real{0.1020}}|%
  >{\raggedright\arraybackslash}p{(\linewidth - 16\tabcolsep - 1pt - 9\arrayrulewidth - 6pt) * \real{0.0888}}|%
  >{\raggedright\arraybackslash}p{(\linewidth - 16\tabcolsep - 1pt - 9\arrayrulewidth - 6pt) * \real{0.1186}}|%
  >{\raggedright\arraybackslash}p{(\linewidth - 16\tabcolsep - 1pt - 9\arrayrulewidth - 6pt) * \real{0.1629}}|%
  >{\raggedright\arraybackslash}p{(\linewidth - 16\tabcolsep - 1pt - 9\arrayrulewidth - 6pt) * \real{0.0889}}|%
  >{\raggedright\arraybackslash}p{(\linewidth - 16\tabcolsep - 1pt - 9\arrayrulewidth - 6pt) * \real{0.1887}}|@{}}
\hline
\begin{minipage}[b]{\hsize}\raggedright
\textbf{Источник}
\end{minipage} & \begin{minipage}[b]{\hsize}\raggedright
\textbf{Спектральный тип шума}
\end{minipage} & \begin{minipage}[b]{\hsize}\raggedright
\textbf{Скорость генерации}
\end{minipage} & \begin{minipage}[b]{\hsize}\raggedright
\textbf{Качество энтропии}
\end{minipage} & \begin{minipage}[b]{\hsize}\raggedright
\textbf{Сложность реализации}
\end{minipage} & \begin{minipage}[b]{\hsize}\raggedright\textbf{Устойчивость}\\\textbf{к внеш.\,воздействиям}%
\end{minipage} & \begin{minipage}[b]{\hsize}\raggedright\textbf{Защита от}\\\textbf{атак / модель.}%
\end{minipage} & \begin{minipage}[b]{\hsize}\raggedright\textbf{Коммент.,}\\\textbf{примеры}%
\end{minipage} \\ \hline
\endhead
\endlastfoot
\textbf{Тепловой шум} & \multirow{2}{\hsize}{Белый шум} & Средняя
(\textasciitilde10⁴--10⁶ бит/с) & Среднее & Низкая & Средняя (зависит от
сопротивления, температуры) & Низкая--средняя & Часто используется в
IC-TRNG, требует постобработки \\ \cline{1-1}\cline{3-8}
\textbf{Дробовой шум} & & Средняя--высокая & Среднее-высокое & Средняя &
Средняя--высокая (дискретность тока более устойчива) & Средняя & Широко
применяется в простых TRNG (на туннельных диодах) \\ \hline
\textbf{Лавинный шум} & Широкополосный с пиками & Высокая
(\textasciitilde10⁶--10⁸ бит/с) & Высокое & Средняя--высокая &
Средняя--высокая (зависит от порога пробоя) & Средняя- высокая &
Популярен в лавинных фотодиодах, хорош для встроенных TRNG \\ \hline
\textbf{Мерцающий шум (1/f)} & Розовый шум & Низкая (\textless10³ бит/с)
& Низкое & Низкая & Низкая (параметры сильно "плавают", нестабильность
спектра) & Низкая & Подвержен внешнему влиянию, требует усиленной
фильтрации \\ \hline
\textbf{Квантово-оптический (фотонный делитель)} & \multirow{2}{\hsize}{Белый
(дискретные события)} & Очень высокая (\textasciitilde10⁶--10⁹ бит/с) &
Очень высокое & Высокая & Высокая (если стабилизирован лазер, фотонный
канал, питание) & Очень высокая & QRNG, QKD, сертифицированные квантовые
генераторы \\ \cline{1-1}\cline{3-8}
\textbf{Радиоактивный распад} & & Очень низкая (\textasciitilde1--10³
бит/с) & Очень высокое & Очень высокая & Очень высокая (независим от
внешней среды) & Очень высокая & Эталонные ГИСП, используются в высокозащищённых системах \\ \hline
\textbf{Механический источник} & \multirow{2}{\hsize}{\raggedright Спектр зависит от
конфигурации (часто псевдобелый)} & Низкая (\textless10² бит/с) & Низкое
& Низкая & \multirow{2}{\hsize}{\raggedright Низкая (вибрации, шумы подвержены массе факторов)} & Очень низкая &
Демонстрационные устройства, устаревшие игровые автоматы \\ \cline{1-1}\cline{3-5}\cline{7-8}
\textbf{Акустический источник} & & Низкая & Низкое & Низкая & & Очень
низкая & Вспомогательная энтропия для ОС (например, Linux entropy
pool) \\ \hline
\end{longtable}
}}
\vspace{0.3\baselineskip}
\raggedright
{\scriptsize\setlength{\itemsep}{0pt}\setlength{\parsep}{0pt}\setlength{\topsep}{0pt}\setlength{\parskip}{0pt}

\textbf{Пояснение.}

\textbf{Качество энтропии:}
\begin{itemize}\setlength{\itemsep}{0pt}\setlength{\parskip}{0pt}
\item Низкое --- сильно коррелированные значения, шум подвержен дрейфу, требуется мощная постобработка;
\item Среднее --- частично предсказуемые биты, возможен bias, возможна частичная корреляция;
\item Высокое --- почти идеальная равномерность и непредсказуемость, даже до постобработки.
\end{itemize}

\textbf{Защита от атак/моделирования:}
\begin{itemize}\setlength{\itemsep}{0pt}\setlength{\parskip}{0pt}
\item Низкая --- источник уязвим к подделке, легко смоделировать или заглушить;
\item Средняя --- требуется физический доступ или точная настройка атаки;
\item Высокая --- крайне сложно предсказать, атака невозможна без вмешательства в систему.
\end{itemize}
}
\end{landscape}
